\documentclass[12pt,a4paper]{article}
\usepackage[a4paper,margin=15mm]{geometry}
\usepackage[T1]{fontenc}
\usepackage{amsmath,amssymb}
\usepackage{enumitem}
\pagestyle{plain}
\setlength{\parindent}{0pt}
\setlength{\parskip}{3pt}
\setlist[enumerate,1]{label=(\alph*),leftmargin=*,itemsep=1pt,topsep=1pt,parsep=0pt}

\begin{document}

\begin{center}
  {\Large\bfseries The Laws of Indices}\\[2pt]
  {\large Answer Key}
\end{center}

\textbf{Q1. Multiplying powers}
\begin{enumerate}
\item \(5^6\)
\item \(x^9\)
\item \(3^7\)
\item \(10^7\)
\item \(7^5\)
\item \(y^{11}\)
\item \(m^6\)
\item \(a^9\)
\item \(7\) \quad (missing index)
\item \(7\) \quad (missing index)
\end{enumerate}

\textbf{Q2. Dividing powers}
\begin{enumerate}
\item \(6^4\)
\item \(x^5\)
\item \(5^6\)
\item \(y^7\)
\item \(9^5\)
\item \(a\)
\item \(3^7\)
\item \(p^5\)
\item \(9\) \quad (missing index)
\item \(9\) \quad (missing index)
\end{enumerate}

\textbf{Q3.} Missing powers: \(3^1,3^{-1},3^{-2}\). Missing values: \(3,1,\frac19,\frac1{27}\).

\textbf{Q4. Zero and negative indices}
\begin{enumerate}
\item \(1\)
\item \(\frac16\)
\item \(\frac19\)
\item \(\frac1{16}\)
\item \(\frac1{100}\)
\item \(\frac1{125}\)
\item \(1\)
\item \(\frac34\)
\end{enumerate}

\textbf{Q5. Using the laws with negative indices}
\begin{enumerate}
\item \(2^2\)
\item \(7^{-3}\)
\item \(5^{-2}\)
\item \(x^{-5}\)
\item \(3^{-6}\)
\item \(a^5\)
\item \(y^2\)
\item \(1\)
\item \(4^{-3}\)
\item \(b^{-4}\)
\end{enumerate}

\textbf{Q6. Powers of powers}
\begin{enumerate}
\item \(5^6\)
\item \(x^{20}\)
\item \(3^8\)
\item \(y^{15}\)
\item \(10^{-6}\)
\item \(a^8\)
\item \(7^{-5}\)
\item \(1\)
\item \(3\) \quad (missing index)
\item \(-2\) \quad (missing index)
\end{enumerate}

\textbf{Q7. Writing numbers as powers of \(2\)}
\begin{enumerate}
\item \(2^6\)
\item \(2^{12}\)
\item \(2^4\)
\item \(2^{-5}\)
\item \(2^{-8}\)
\item \(2^1=2\)
\end{enumerate}

\textbf{Q8. Simplifying expressions with indices}
\begin{enumerate}
\item \(10x^7\)
\item \(3p^3\)
\item \(12y^7\)
\item \(4m^6\)
\item \(12a^3b^4\)
\item \(5x^3y^3\)
\item \(9x^4\)
\item \(8a^9b^3\)
\item \(12x^3\)
\item \(5y^{-2}=\frac{5}{y^2}\)
\item \(25k^{-2}=\frac{25}{k^2}\)
\end{enumerate}

\textbf{Q9. Area and volume with indices}
\begin{enumerate}
\item \(6x^7\ \text{cm}^2\)
\item \(8a^9\ \text{cm}^3\)
\item \(6x^5\ \text{cm}\)
\end{enumerate}

\textbf{Q10.} Examples:
\[
4x^5\times4x^7=16x^{12},\qquad
32x^{15}\div2x^3=16x^{12},\qquad
(4x^6)^2=16x^{12}.
\]

\textbf{Q11. Fractions as the base}
\begin{enumerate}
\item \(\frac9{16}\)
\item \(\frac{25}{4}\)
\item \(\frac18\)
\item \(\frac8{27}\)
\item \(4\)
\item \(\frac43\)
\item \(9\)
\item \(\frac{27}{8}\)
\item \(\frac4{25}\)
\item \(1\)
\end{enumerate}

\textbf{Q12. Decimals as bases}
\begin{enumerate}
\item \(0.6^{-2}=\frac{25}{9}\)
\item \(0.25^{-3}=64\)
\item \(0.3^2=\frac9{100}=0.09\)
\item \(0.2^3=\frac1{125}=0.008\)
\item \(0.5^{-1}=2\)
\item \(0.5^{-3}=8\)
\item \(0.1^{-2}=100\)
\item \(0.2^{-2}=25\)
\item \(1.5^2=\frac94=2.25\)
\item \(0.75^{-2}=\frac{16}{9}\)
\item \(2.5^{-2}=\frac4{25}=0.16\)
\item \(0.8^{-1}=\frac54=1.25\)
\end{enumerate}

\textbf{Q13. Using the laws, then working out the value}
\begin{enumerate}
\item \(4\)
\item \(\frac94\)
\item \(25\)
\item \(1000\)
\item \(\frac{25}{4}\)
\item \(16\)
\end{enumerate}

\textbf{Q14. Tiles and cubes}
\begin{enumerate}
\item \(16\) tiles; this agrees with \(0.25^{-2}=16\).
\item \(125\) small cubes.
\end{enumerate}

\textbf{Q15. Reasoning questions}
\begin{enumerate}
\item Sam is not correct: \(0.5^{-2}=4\), which is positive.
\item
\[
4^{-1}=\frac14
\;<\;
\left(\frac32\right)^{-2}=\frac49
\;<\;
0.5^{-3}=8
\;<\;
0.1^{-1}=10
\;<\;
0.2^{-2}=25.
\]
\item Many answers are possible, for example:
\[
16=4^2=16^1=0.25^{-2}=\left(\frac14\right)^{-2}
=\left(\frac1{16}\right)^{-1}=0.0625^{-1}
=(-4)^2=(-2)^4=2^4=0.5^{-4}.
\]
\end{enumerate}

\end{document}